\documentclass[11pt]{article}
\usepackage[T1]{fontenc}
\usepackage{lmodern,amsmath,amssymb,amsthm,geometry,microtype,xcolor}
\geometry{margin=1in}
\usepackage[colorlinks=true,allcolors=blue!45!black]{hyperref}
\newtheorem{theorem}{Theorem}
\title{A finite counterexample to the unrestricted\newline additive-count statement in MathDB}
\author{Myeongjin Kang}
\date{October 26, 2026}
\begin{document}
\maketitle
\begin{abstract}
MathDB question 334772 asserts that, for an arbitrary set of inputs and an
arbitrary family of arithmetic functions, the number satisfying an additive
relation cannot increase with the number of summands. A singleton input set
and a family of two real-valued functions give counts one and two at arities
two and three. The complete counterexample, including the universal additive
conditions and the cardinalities of the filtered function family, is proved
in Lean. This is a correction to the literal corpus statement. The cited
arXiv manuscript does not contain that unrestricted conjecture.
\end{abstract}
\noindent\textit{Scope and credit.} This elementary observation is new to
this project; no worldwide priority or substantive open-problem resolution
is claimed. In particular, the refuted assertion is not attributed to Park.
This paper is AI-assisted. Original prose is CC BY 4.0 and original code MIT.

\section{The precise target and its source limitation}
For $E\subseteq\mathbb N$, a finite family $S$ of real-valued arithmetic
functions, and a positive integer $k$, define
\[
 \mathcal A_k(E)=\left\{f:\mathbb N\to\mathbb R:
 f\left(\sum_{i=1}^k a_i\right)=\sum_{i=1}^k f(a_i)
 \text{ for every }(a_1,\ldots,a_k)\in E^k\right\},
\]\[
 n(E,S,k)=|S\cap\mathcal A_k(E)|.
\]
MathDB question 334772 states a general monotonicity conjecture for these
counts: $n(E,S,k)\ge n(E,S,\ell)$ when $k\le\ell$~\cite{MathDB}.
We use the ordinary meaning of counting functions satisfying the displayed
relation. The statement places no multiplicativity, infinitude, or closure
assumption on $S$ or $E$. Our example has positive inputs only, so the
convention about including zero in $\mathbb N$ has no effect.

The record cites Park's arXiv:1612.00897~\cite{Park}. Its sole arXiv version,
v1 of December 2, 2016, has been read in full. Its main theorem concerns
multiplicative functions additive on sums of positive squares for $k\ge3$.
It does not state the unrestricted count conjecture or define $n(E,S,k)$.
Consequently the target here is the MathDB wording itself, and its cited
attribution is not established. We make no claim about any conjecture of
Park or about a version with further restrictions.

\newpage
\section{A complete finite counterexample}
\begin{theorem}
Let $E=\{1\}$ and $S=\{u,v\}$, where
\[
 u(n)=n,\qquad
 v(n)=\begin{cases}0&n=2,\\n&n\ne2.\end{cases}
\]
Then $n(E,S,2)=1$ and $n(E,S,3)=2$. In particular, the unrestricted
monotonicity assertion is false.
\end{theorem}
\begin{proof}
The functions are distinct because $u(2)=2$ and $v(2)=0$. The identity
function $u$ satisfies the additive relation for every $E$ and $k$.
For $E=\{1\}$ there is exactly one ordered $k$-tuple of inputs: every
entry is one. The relation for $v$ is therefore equivalent to $v(k)=kv(1)$.
At $k=2$, its left side is zero and its right side is two, so it fails.
At $k=3$, both sides are three, so it holds. Thus precisely one of the two
functions satisfies the two-summand relation, whereas both satisfy the
three-summand relation. Since $2\le3$ and $1<2$, the claimed inequality
has the wrong direction for this example.
\end{proof}

\section{Formal coverage and reproducibility}
The module \texttt{AdditiveMonotonicity.lean} defines real-valued functions
on natural numbers, the universal predicate \texttt{KAdditive}, and
\texttt{count} as the actual cardinality of a filtered finite set of
functions. The functions are not represented merely by labels. Function
extensionality reduces every tuple in the singleton input set to the
constant tuple. Separate declarations establish failure at arity two,
validity at arity three, distinctness of the functions, and both exact
cardinalities. The final theorem
\texttt{AdditiveMonotonicity.monotonicity\_counterexample} combines
$2\le3$ with the strict increase of the actual counts.

All claims of the theorem are covered by Lean 4.34.1 using the pinned
mathlib dependencies in the accompanying manifest. The isolated verifier
compiles the entire module and checks every declaration's transitive
axioms; it permits only Lean's standard logical axioms. The source audit,
independent review, paper, and reproduction instructions accompany the
certificate. Only this single source question is counted as an outcome.

\begin{thebibliography}{9}
\bibitem{MathDB} MathDB, \emph{Monotonicity of the number of functions
determined by additive relations}, problem 334772, accessed October 26,
2026. \url{https://mathdb.com/p/334772}.
\bibitem{Park} P.-S. Park, \emph{$k$-additive uniqueness of the set of squares
for multiplicative functions}, arXiv:1612.00897v1, December 2, 2016.
\url{https://arxiv.org/abs/1612.00897v1}.
\end{thebibliography}
\end{document}
